% !TEX root = ../main.tex

\section{One-Dimensional Pedagogical Model}\label{sec:1d-model}

\subsection{Model Definition}\label{subsec:model-def}

We consider a one-dimensional lattice populated with two species of particles:
\begin{itemize}
    \item $(+)$ particles: move rightward to $i+1$ at a constant rate $\lambda$
    \item $(-)$ particles: move leftward to $i-1$ at a constant rate $\lambda$
\end{itemize}

Let $n_i^\pm$ denote the number of species $\pm$ particles at lattice site $i$. The dynamics of the system consists of two processes:

\begin{keybox}[Model Dynamics]
\begin{enumerate}
    \item \textbf{Streaming}: Particles move from $i$ to $i\pm 1$ at rate $\lambda$
    \item \textbf{Conversion}: Particles convert from $+$ to $-$ (or vice versa) at rate $r_\pm(n_i^+, n_i^-)$
\end{enumerate}
\end{keybox}

The choice of conversion rates reflects the physics of antialigning interactions. In this paper, the simplest nonlinear form is chosen (ignoring self-interactions):
\begin{equation}\label{eq:conversion-rate}
\boxed{r_\pm(n^+, n^-) = r_0 + (n^\pm - 1)r_1}
\end{equation}

where $r_0$ is the ``spontaneous conversion rate'' and $r_1$ is the antialigning interaction strength associated with the local magnetization.

\begin{remark}[Note on the form of conversion rates]
Here we require $r_1$ to be much larger than $r_0$, because we are focusing on conversions induced by local magnetization. Keeping only terms up to $O(n^+, n^-)$ in the conversion rate is crucial---higher-order terms would lead to derivatives higher than second order in the Fokker-Planck equation, which cannot be solved directly via the corresponding Langevin equation.
\end{remark}

\subsection{Master Equation}\label{subsec:master-eq}

Let $P(\{n_i^+, n_i^-\})$ be the probability distribution of finding a specific set of occupation numbers. The master equation reads:
\begin{equation}\label{eq:master-eq}
\partial_t P = -S^\lambda_{\text{out}} + S^\lambda_{\text{in}} - C^r_{\text{out}} + C^r_{\text{in}}
\end{equation}

The individual terms represent:
\begin{itemize}
    \item $S^\lambda_{\text{out}}$: Probability outflow due to streaming
    \item $S^\lambda_{\text{in}}$: Probability inflow due to streaming
    \item $C^r_{\text{out}}$: Probability outflow due to conversion
    \item $C^r_{\text{in}}$: Probability inflow due to conversion
\end{itemize}

\subsubsection{Explicit Expressions for Probability Currents}

\textbf{Streaming terms} (outflow):
\begin{align}
S^\lambda_{\text{out}} &= \sum_i \left[\lambda n_i^+ + \lambda n_i^-\right] P \label{eq:stream-out}\\
C^r_{\text{out}} &= \sum_i \left[r_+ n_i^+ + r_- n_i^-\right] P \label{eq:conv-out}
\end{align}

\textbf{Streaming terms} (inflow, contributions from left and right neighbors):
\begin{align}
S^\lambda_{\text{in}} &= \sum_i \Big[\lambda(n_{i-1}^+ + 1)P(\ldots, n_{i-1}^+ + 1, n_i^+ - 1, \ldots) \nonumber\\
&\qquad + \lambda(n_{i+1}^- + 1)P(\ldots, n_i^- - 1, n_{i+1}^- + 1, \ldots)\Big] \label{eq:stream-in}
\end{align}

\textbf{Conversion terms} (inflow):
\begin{align}
C^r_{\text{in}} &= \sum_i \Big[r_+(n_i^+ + 1, n_i^- - 1)(n_i^+ + 1)P(\ldots, n_i^+ + 1, n_i^- - 1, \ldots) \nonumber\\
&\qquad + r_-(n_i^+ - 1, n_i^- + 1)(n_i^- + 1)P(\ldots, n_i^+ - 1, n_i^- + 1, \ldots)\Big] \label{eq:conv-in}
\end{align}

\subsection{Poisson Representation}\label{subsec:poisson-rep}

Directly solving the master equation is infeasible because for interacting systems (i.e., nontrivial conversion rates), the right-hand side is nonlinear. Therefore, we adopt the \textbf{Poisson representation}\cite{gardiner1977}, a method that transforms a discrete stochastic process into a continuous one.

\subsubsection{Core Idea}

\begin{notebox}[Core Idea of the Poisson Representation]
Replace the discrete occupation numbers $n_i^\pm$ (a stochastic process on positive integers) with continuous Poisson rates $\alpha_i, \beta_i$ (corresponding to $+$ and $-$ particles, respectively). These rates serve as parameters of Poisson processes that then generate the occupation numbers $n_i^\pm$.

Statistical moments are obtained by averaging over realizations of the Poisson process for fixed fields, as well as over the distribution of these fields.
\end{notebox}

The relationship between the probability distribution $P(\{n_i^+, n_i^-\})$ and the quasiprobability distribution $f(\{\alpha_i, \beta_i\})$ is:
\begin{equation}\label{eq:poisson-rep}
\boxed{P = \int \prod_i \dd\alpha_i \, \dd\beta_i \; e^{-\alpha_i}\frac{\alpha_i^{n_i^+}}{n_i^+!} \cdot e^{-\beta_i}\frac{\beta_i^{n_i^-}}{n_i^-!} \; f}
\end{equation}

\subsubsection{Mismatch Between General Discrete Processes and Poisson Processes}

Due to a ``mismatch'' between the generality of discrete processes and local Poisson processes, a higher level of abstraction is required:
\begin{enumerate}
    \item The local rates $\alpha_i, \beta_i$ are \textbf{complex numbers}
    \item The distribution $f$ is actually a \textbf{quasidistribution}, as it can also be locally negative
\end{enumerate}

These extensions are not unphysical, because the Poisson fields are auxiliary mathematical objects that are not directly observable. \textbf{Observable} are the moments of the physical process, which are positive and real.

\begin{remark}[Intuitive understanding of imaginary numbers]
The contribution of imaginary numbers to the Poisson field can be viewed as a manifestation of deviations from Poisson statistics, which explicitly couples the mean and variance of the random variables.

Compare a Poisson variable $n$ with rate $\lambda$ and a variable $m$ with rate $\lambda + i\xi$, where $\xi$ is a zero-mean random variable with $\langle \xi^2 \rangle = c$. For a Poisson process, we have:
\[
\langle n \rangle = \langle n^2 \rangle_c = \lambda
\]
But for $m$:
\[
\langle m \rangle = \lambda, \quad \langle m^2 \rangle_c = \lambda - c^2
\]
Imaginary noise leads to \textbf{narrowing} of the microscopic distribution, while real noise leads to \textbf{widening}.
\end{remark}

\subsubsection{Replacement Rules: From Master Equation to Fokker-Planck Equation}

By substituting the Poisson representation into the master equation and using integration by parts, we obtain the following replacement rules. Below, arguments in $f$ and $P$ that are consistent with the aforementioned default parameters are omitted.

\begin{derivation}[Replacement rules for linear terms (streaming and spontaneous conversion)]
\textbf{Outflow term} (rightward motion):
\begin{equation}\label{eq:replace-out-right}
n_i^+ P \;\to\; (\alpha_i f - \partial_{\alpha_i}\alpha_i f)
\end{equation}

\textbf{Outflow term} (leftward motion):
\begin{equation}\label{eq:replace-out-left}
n_i^- P \;\to\; (\beta_i f - \partial_{\beta_i}\beta_i f)
\end{equation}

\textbf{Inflow term} (from the left):
\begin{equation}\label{eq:replace-in-left}
(n_{i-1}^+ + 1)P(\ldots, n_{i-1}^+ + 1, n_i^+ - 1, \ldots) \;\to\; (\alpha_{i-1} - \partial_{\alpha_i}\alpha_{i-1})f
\end{equation}

\textbf{Inflow term} (from the right):
\begin{equation}\label{eq:replace-in-right}
(n_{i+1}^- + 1)P(\ldots, n_i^- - 1, n_{i+1}^- + 1, \ldots) \;\to\; (\beta_{i+1} - \partial_{\beta_i}\beta_{i+1})f
\end{equation}
\end{derivation}

\begin{derivation}[Replacement rules for nonlinear terms (conversion process)]
\textbf{Conversion from $(+)$ to $(-)$}:
\begin{equation}\label{eq:replace-conv-plus}
n_i^+(n_i^+ - 1)P \;\to\; (1 - \partial_{\alpha_i})^2 \alpha_i^2 f
\end{equation}

\textbf{Conversion from $(-)$ to $(+)$}:
\begin{equation}\label{eq:replace-conv-minus}
n_i^-(n_i^- - 1)P \;\to\; (1 - \partial_{\beta_i})^2 \beta_i^2 f
\end{equation}

\textbf{Conversion of $(+)$ particles induced by $(-)$ particles}:
\begin{equation}\label{eq:replace-conv-mix-plus}
(n_i^- + 1)P(\ldots, n_i^+ - 1, n_i^- + 1, \ldots) \;\to\; (-1 + \partial_{\alpha_i})(-1 + \partial_{\beta_i})\alpha_i^2 f
\end{equation}

\textbf{Conversion of $(-)$ particles induced by $(+)$ particles}:
\begin{equation}\label{eq:replace-conv-mix-minus}
(n_i^+ + 1)P(\ldots, n_i^+ + 1, n_i^- - 1, \ldots) \;\to\; (-1 + \partial_{\alpha_i})(-1 + \partial_{\beta_i})\beta_i^2 f
\end{equation}
\end{derivation}

\subsubsection{Detailed Derivation: Obtaining the Replacement Rules}

Taking Eq.~(\ref{eq:replace-out-right}) as an example, we show the detailed derivation:
\begin{align}
n_i^+ P &= \int \prod_j \dd\alpha_j \, \dd\beta_j \; n_i^+ \frac{e^{-\alpha_i}\alpha_i^{n_i^+}}{n_i^+!} \cdot \frac{e^{-\beta_i}\beta_i^{n_i^-}}{n_i^-!} \; f \nonumber\\
&= \int \prod_j \dd\alpha_j \, \dd\beta_j \; \frac{e^{-\alpha_i}\alpha_i^{n_i^+}}{(n_i^+ - 1)!} \cdot \frac{e^{-\beta_i}\beta_i^{n_i^-}}{n_i^-!} \; f \nonumber\\
&= \int \prod_j \dd\alpha_j \, \dd\beta_j \; \frac{e^{-\alpha_i}\alpha_i^{n_i^+}}{n_i^+!} \cdot \frac{e^{-\beta_i}\beta_i^{n_i^-}}{n_i^-!} \; \alpha_i f \nonumber\\
&\quad - \int \prod_j \dd\alpha_j \, \dd\beta_j \; (-\partial_{\alpha_i})\left[e^{-\alpha_i}\frac{\alpha_i^{n_i^+}}{n_i^+!}\right] \cdot \frac{e^{-\beta_i}\beta_i^{n_i^-}}{n_i^-!} \; \alpha_i f \nonumber\\
&= \text{form of } P \text{ after } n_i^+ \text{ is replaced} \nonumber
\end{align}

After integration by parts (assuming boundary terms vanish), the replacement rule is obtained.

For nonlinear terms, e.g., Eq.~(\ref{eq:replace-conv-plus}):
\begin{align}
n_i^+(n_i^+ - 1)P &= \int \dd\alpha_i \; n_i^+(n_i^+ - 1) \frac{e^{-\alpha_i}\alpha_i^{n_i^+}}{n_i^+!} \; f \nonumber\\
&= \int \dd\alpha_i \; \frac{e^{-\alpha_i}\alpha_i^{n_i^+}}{(n_i^+ - 2)!} \; f \nonumber\\
&= \int \dd\alpha_i \; (1 - \partial_{\alpha_i})^2 \left[\frac{e^{-\alpha_i}\alpha_i^{n_i^+}}{n_i^+!}\right] \alpha_i^2 f \nonumber\\
&\to\; (1 - \partial_{\alpha_i})^2 \alpha_i^2 f \nonumber
\end{align}

\subsection{Fokker-Planck Equation}\label{subsec:fokker-planck}

Collecting the above replacement rules, we obtain the Fokker-Planck equation for the quasiprobability density $f(\{\alpha_i, \beta_i\})$:
\begin{equation}\label{eq:fpe}
\boxed{\partial_t f = F_\lambda + F_r + J}
\end{equation}

where the individual terms are:

\subsubsection{Streaming Contribution $F_\lambda$}
\begin{align}
F_\lambda &= \lambda \sum_i \Big[(\alpha_{i-1} - \alpha_i) + \partial_{\alpha_i}(\alpha_i - \alpha_{i-1})\Big]f \nonumber\\
&\quad - \lambda \sum_i \Big[(\beta_i - \beta_{i+1}) + \partial_{\beta_i}(\beta_{i+1} - \beta_i)\Big]f \label{eq:f-lambda}
\end{align}

\subsubsection{Conversion Contribution $F_r$}
\begin{align}
F_r &= r_0 \sum_i \Big[\partial_{\alpha_i}(\alpha_i - \beta_i) - \partial_{\beta_i}(\alpha_i - \beta_i)\Big]f \nonumber\\
&\quad + r_1 \sum_i \Big[\partial_{\alpha_i} - \partial_{\beta_i}\Big](\alpha_i + \beta_i)(\alpha_i - \beta_i)f \label{eq:f-r}
\end{align}

\subsubsection{Noise Term $J$}

\begin{equation}\label{eq:noise-term}
\boxed{J = -r_1 \sum_i \Big[\partial_{\alpha_i}^2 \alpha_i^2 + \partial_{\beta_i}^2 \beta_i^2\Big]f}
\end{equation}

The last line contains second-derivative terms, which manifest as \textbf{noise} in the corresponding Langevin equation. Note that the coefficient of this noise term is \textbf{negative}, which ultimately leads to the emergence of imaginary noise.

\subsection{Deterministic Equations of Motion}\label{subsec:det-eom}

From the Fokker-Planck equation, we can read off the deterministic equations (ignoring the noise term). Changing the notation from $\lambda$ to $v$ (for consistency with tradition), we obtain:
\begin{align}
\partial_t \alpha\big|_{\text{det}} &= -v\nabla\alpha - r_0(\alpha - \beta) - r_1(\alpha + \beta)(\alpha - \beta) \label{eq:det-alpha}\\
\partial_t \beta\big|_{\text{det}} &= v\nabla\beta + r_0(\alpha - \beta) + r_1(\alpha + \beta)(\alpha - \beta) \label{eq:det-beta}
\end{align}

Note that here $\nabla$ is shorthand for nearest-neighbor differences, and the \textbf{continuous limit has not yet been taken}.

\subsection{Noise Analysis}\label{subsec:noise-analysis}

To determine the noise term, we write the Fokker-Planck equation in standard form:
\begin{equation}\label{eq:fpe-standard}
\partial_t f = -\sum_i \partial_{x_i}(\mu_i f) + \sum_{i,j} \partial_{x_i}\partial_{x_j}[D_{ij} f]
\end{equation}

The diffusion matrix $D_{ij}$ is:
\begin{equation}\label{eq:diffusion-matrix}
D_{ij} = -r_1 \begin{pmatrix} \alpha_i^2 & 0 \\ 0 & \beta_i^2 \end{pmatrix}
\end{equation}

Introducing the variable transformation:
\begin{equation}\label{eq:variable-change}
\tilde{\rho}_i = \alpha_i + \beta_i \quad (\text{total density}), \qquad \tilde{m}_i = \alpha_i - \beta_i \quad (\text{magnetization})
\end{equation}

The diffusion matrix becomes:
\begin{equation}\label{eq:diffusion-rho-m}
D_{ij} = -r_1 \begin{pmatrix} \frac{\tilde{\rho}^2 + \tilde{m}^2}{4} & \frac{-2\tilde{\rho}\tilde{m}}{4} \\ \frac{-2\tilde{\rho}\tilde{m}}{4} & \frac{\tilde{\rho}^2 + \tilde{m}^2}{4} \end{pmatrix}
\end{equation}

\subsubsection{Matrix Square Root and Imaginary Noise}

Translating the Fokker-Planck equation into Langevin equations requires computing the matrix square root $B = \sqrt{2D}$. In the physically relevant limit $\tilde{\rho} \approx \rho_0$, $\tilde{m} \approx 0$, to zeroth order in $\tilde{m}$:
\begin{equation}\label{eq:matrix-sqrt}
B = i\rho_0\sqrt{\frac{r_1}{2}} \begin{pmatrix} 1 & -1 \\ -1 & 1 \end{pmatrix}
\end{equation}

Here the factor $i = \sqrt{-1}$ appears! This is the hallmark result of the Poisson representation.

\begin{notebox}[Physical Meaning of Imaginary Noise]
The two eigenvalues are:
\[
\lambda_\rho = 0, \qquad \lambda_m = i\sqrt{2r_1\rho_0}
\]
$\lambda_\rho = 0$ means that \textbf{the density field has no noise}, while the magnetization field has imaginary noise.

This does not imply that any observable is imaginary. The appearance of imaginary noise signals a deviation from Poisson statistics---\textbf{in the Poisson representation, only complex fields can describe non-Poissonian statistics}, because the mean and variance of a Poisson distribution must be equal, and to break this equality, the imaginary part of the field must be nonzero.
\end{notebox}

\subsection{Langevin Equations}\label{subsec:langevin}

Summarizing the above analysis, the coupled Langevin equations for the density and magnetization fields are:

\begin{eqbox}[Langevin Equations]
\begin{align}
\partial_t \tilde{\rho} &= -v\nabla\tilde{m} \label{eq:langevin-rho}\\
\partial_t \tilde{m} &= -v\nabla\tilde{\rho} - r_0\tilde{m} - r_1\tilde{\rho}\tilde{m} + i\sqrt{2r_1\tilde{\rho}}\,\xi \label{eq:langevin-m}
\end{align}
\end{eqbox}

where $\xi$ is Gaussian uncorrelated noise satisfying:
\begin{equation}\label{eq:noise-correlation}
\langle \xi \rangle = 0, \qquad \langle \xi(t)\xi(t') \rangle = \delta(t - t')
\end{equation}

\begin{remark}[Term-by-term physical interpretation]
\textbf{Density equation (\ref{eq:langevin-rho}):}
\begin{itemize}
    \item $\partial_t \tilde{\rho}$: Local time rate of change of the total density field.
    \item $-v\nabla\tilde{m}$: \textbf{Density current driven by magnetization gradient}. Since $(+)$ particles move rightward and $(-)$ particles move leftward, if there are more rightward-moving particles to the right than to the left at a given location (i.e., $\nabla\tilde{m}>0$), the density at that location decreases, and vice versa. The negative sign reflects the net transport of particles from regions of high magnetization to regions of low magnetization.
\end{itemize}

\textbf{Magnetization equation (\ref{eq:langevin-m}):}
\begin{itemize}
    \item $\partial_t \tilde{m}$: Time rate of change of the magnetization (i.e., the difference between the number of rightward and leftward particles).
    \item $-v\nabla\tilde{\rho}$: \textbf{Magnetization current driven by density gradient}. If the density increases to the right ($\nabla\tilde{\rho}>0$), more rightward-moving particles flow in from the left, causing the local magnetization to decrease.
    \item $-r_0\tilde{m}$: \textbf{Magnetization relaxation due to spontaneous conversion}. $r_0$ is the spontaneous flipping rate, which causes $(+)$ and $(-)$ particles to convert into each other, driving the system toward an unpolarized state ($\tilde{m}=0$); this is an exponential decay term.
    \item $-r_1\tilde{\rho}\tilde{m}$: \textbf{Nonlinear relaxation due to antialigning interaction}. This is the most crucial term in this paper: the stronger the local magnetization, the faster particles flip to the opposite direction due to the ``antialigning'' mechanism; moreover, this effect is proportional to the local total density (the higher the density, the more frequent the collisions/interactions). This term is the key to suppressing magnetization and, in turn, suppressing density fluctuations.
    \item $i\sqrt{2r_1\tilde{\rho}}\,\xi$: \textbf{Imaginary noise term}. This is the hallmark result of the Poisson representation. The discreteness and nonlinearity of the conversion process lead to intrinsic stochastic fluctuations in the magnetization dynamics; the negative definiteness of the diffusion matrix $D$ makes the noise coefficient purely imaginary, meaning that the system exhibits \textbf{sub-Poissonian} statistics (fluctuations narrower than the Poisson distribution). Note that there is no noise at all in the density equation, which is an important signature of the strong suppression of fluctuations in this model.
\end{itemize}
\end{remark}

\begin{remark}[Connection to existing literature]
The deterministic version of this set of equations (setting $r_0 \equiv 0$) has the same structure as the closed hydrodynamics of antialigning particles\cite{boltz2025} in one dimension. General Toner-Tu-type approaches\cite{toner1998} also encompass these dynamics, but do not make specific predictions for $r_0$ and $r_1$.
\end{remark}

\subsection{Analytical Calculation of the Structure Factor}\label{subsec:structure-factor}

\subsubsection{Linearization and Simplification}

Consider small perturbations around the uniform unpolarized state:
\begin{equation}\label{eq:linearization}
\tilde{\rho} = \rho_0 + \delta\tilde{\rho}, \qquad \tilde{m} \text{ treated as a small quantity}
\end{equation}

Approximating the multiplicative noise as additive noise (ignoring the $\delta\tilde{\rho}\,\xi$ term), we obtain a closed second-order equation:
\begin{equation}\label{eq:second-order-eq}
\partial_t^2 \delta\tilde{\rho} = v^2\nabla^2\delta\tilde{\rho} - (r_0 + r_1\rho_0)\partial_t\delta\tilde{\rho} - i\sqrt{2r_1\rho_0}\,v\nabla\xi
\end{equation}

\subsubsection{Overdamped Approximation}

\begin{derivation}
Write the \textbf{deterministic} part of the second-order equation~(\ref{eq:second-order-eq}) (ignoring noise for now) in Fourier space, setting $\nabla^2 \to -k^2$:
\begin{equation}\label{eq:damped-oscillator}
\partial_t^2 \delta\tilde{\rho}_k + \gamma\,\partial_t \delta\tilde{\rho}_k + \omega_0^2\,\delta\tilde{\rho}_k = 0
\end{equation}

\begin{notebox}[Notation: Meaning of $\delta\tilde{\rho}_k$]
Recall Eq.~(\ref{eq:linearization}), where we wrote the total density as the sum of a uniform background and a small perturbation: $\tilde{\rho} = \rho_0 + \delta\tilde{\rho}$. The subscript $k$ denotes the \textbf{wavenumber component} after Fourier transforming the spatial coordinate:
\[
\delta\tilde{\rho}_k(t) = \int \dd x\, e^{-ikx}\,\delta\tilde{\rho}(x,t)
\quad\text{(continuous)}\qquad\text{or}\qquad
\delta\tilde{\rho}_k(t) = \sum_j e^{-ikj}\,\delta\tilde{\rho}_j(t)
\quad\text{(discrete lattice)}
\]
Thus $\delta\tilde{\rho}_k$ is the \textbf{Fourier mode amplitude} of the density fluctuation at wavenumber $k$. When taking time derivatives, the spatial derivative $\nabla$ is replaced by $ik$ (or $\nabla^2 \to -k^2$), thereby reducing the partial differential equation to an ordinary differential equation for each $k$.
\end{notebox}
This is the standard form of a \textbf{damped harmonic oscillator}, where:
\begin{itemize}
    \item $\gamma = r_0 + r_1\rho_0$ \quad(damping coefficient, from spontaneous conversion and antialigning interaction)
    \item $\omega_0^2 = v^2 k^2$ \quad(square of the natural frequency, from the directed motion of particles)
\end{itemize}

The characteristic equation of the damped harmonic oscillator is $\lambda^2 + \gamma\lambda + \omega_0^2 = 0$, yielding two eigenvalues:
\begin{equation}\label{eq:eigenvalues-full}
\lambda_{\pm} = \frac{-\gamma \pm \sqrt{\gamma^2 - 4\omega_0^2}}{2}
\end{equation}

\textbf{Overdamped condition}: When $\gamma^2 \gg 4\omega_0^2$, the quantity under the square root is positive, and both eigenvalues are \textbf{real and negative}; the system does not oscillate but decays directly. In this model:
\[
\gamma^2 = (r_0 + r_1\rho_0)^2, \qquad 4\omega_0^2 = 4v^2 k^2
\]
Since $\gamma$ is a constant while $\omega_0 \propto k$, in the \textbf{large-scale limit $k \to 0$}, the condition is automatically satisfied! This is the fundamental reason why the system is overdamped at large scales: the rapid conversion between particles ($\gamma$) smooths out the oscillatory tendency induced by streaming.

\begin{remark}[Why expand the square root? On what approximation is this based?]
Without expansion, the eigenvalues~(\ref{eq:eigenvalues-full}) are wrapped inside a square root, and we \textbf{cannot directly read off} the scaling behavior of the two modes. The core of the overdamped limit is $\gamma^2 \gg 4\omega_0^2$, i.e.,
\[
\epsilon \equiv \frac{4\omega_0^2}{\gamma^2} = \frac{4v^2 k^2}{(r_0 + r_1\rho_0)^2} \ll 1
\]
This is a \textbf{small parameter}. Rewrite the square root as
\[
\sqrt{\gamma^2 - 4\omega_0^2} = \gamma\sqrt{1 - \epsilon}
\]
and Taylor expand $\sqrt{1 - \epsilon}$ (to first order):
\[
\sqrt{1 - \epsilon} = 1 - \frac{\epsilon}{2} + O(\epsilon^2) \approx 1 - \frac{2\omega_0^2}{\gamma^2}
\]
In this way the square root is \textbf{linearized}, and upon substitution the $k$ dependence of the two eigenvalues becomes transparent: one is proportional to $-\gamma$ (fast mode), and the other is proportional to $-k^2$ (slow mode). If the square root were kept, subsequent scaling analysis of the structure factor and diffusion coefficient would be obstructed.
\end{remark}

Expanding the square root in the $k \to 0$ limit:
\[
\sqrt{\gamma^2 - 4\omega_0^2} \approx \gamma\left(1 - \frac{2\omega_0^2}{\gamma^2}\right) = \gamma - \frac{2\omega_0^2}{\gamma}
\]
Substituting into Eq.~(\ref{eq:eigenvalues-full}), we obtain two modes:
\begin{align}
    \lambda_+ &= \frac{-\gamma + (\gamma - 2\omega_0^2/\gamma)}{2} = -\frac{\omega_0^2}{\gamma} = -\frac{v^2 k^2}{r_0 + r_1\rho_0} \label{eq:slow-mode}\\
    \lambda_- &= \frac{-\gamma - (\gamma - 2\omega_0^2/\gamma)}{2} = -\gamma + \frac{\omega_0^2}{\gamma} \approx -(r_0 + r_1\rho_0) \label{eq:fast-mode}
\end{align}

\begin{itemize}
    \item $\lambda_-$ is the \textbf{fast mode} ($\sim -\gamma$), decaying to zero in extremely short times via $e^{\lambda_- t}$.
    \item $\lambda_+$ is the \textbf{slow mode} ($\sim -k^2$), determining the long-time behavior of the system.
\end{itemize}

Therefore, the long-time behavior is dominated by the slow mode, i.e.:
\begin{equation}\label{eq:eigenvalue}
\boxed{\mu = \lambda_+ = -\frac{v^2 k^2}{r_0 + r_1\rho_0}}
\end{equation}
\end{derivation}

This gives the diffusive scaling $\sim e^{-Dk^2t}$, where the diffusion coefficient is:
\begin{equation}\label{eq:diffusivity}
\boxed{D = \frac{v^2}{r_0 + r_1\rho_0}}
\end{equation}

\begin{remark}[Another intuitive understanding of overdamping]
One can also see the essence of overdamping directly from the second-order equation~(\ref{eq:second-order-eq}): the coefficient of the damping term $\gamma\partial_t\delta\tilde{\rho}$ is $\gamma = r_0 + r_1\rho_0$, which at $k \to 0$ is much larger than the effective coefficient of the inertial term $\partial_t^2\delta\tilde{\rho}$ (which can be taken as 1). Therefore, we can \textbf{adiabatically neglect the inertial term} $\partial_t^2\delta\tilde{\rho} \approx 0$, immediately obtaining:
\[
0 \approx -v^2 k^2 \delta\tilde{\rho}_k - (r_0 + r_1\rho_0)\partial_t\delta\tilde{\rho}_k
\quad\Longrightarrow\quad
\partial_t\delta\tilde{\rho}_k = -\frac{v^2 k^2}{r_0 + r_1\rho_0}\delta\tilde{\rho}_k
\]
This is completely consistent with Eq.~(\ref{eq:eigenvalue}). Physically, this means that the magnetization (fast variable) is \textbf{slaved} by the density (slow variable), and the system relaxes by diffusion without oscillation.
\end{remark}

\subsubsection{First-Order Langevin Equation}

Since the deterministic dynamics is overdamped at large scales, we can neglect the inertial term and reduce everything to a first-order Langevin equation (or noisy diffusion equation):
\begin{equation}\label{eq:first-order-langevin}
\partial_t \delta\tilde{\rho}_k = -Dk^2\delta\tilde{\rho}_k + ic\eta_k
\end{equation}

where:
\begin{equation}\label{eq:noise-strength}
c = v\sqrt{\frac{2r_1\rho_0}{r_0 + r_1\rho_0}}
\end{equation}

\subsubsection{Definition of Correlated Noise}

To correctly understand the effective noise statistics in Fourier space, note that $\nabla\xi$ is actually $\xi_{i+1} - \xi_i$. For this purpose, we introduce the zero-drift correlated noise $\eta_i = \eta(i)$:
\begin{equation}\label{eq:correlated-noise}
\langle \eta_i(t)\eta_j(t') \rangle = \delta(t - t') \times \begin{cases} 2, & i = j \\ -1, & i = j \pm 1 \\ 0, & \text{otherwise} \end{cases}
\end{equation}

In Fourier space:
\begin{equation}\label{eq:noise-fourier}
\langle \eta_k(t)\eta_{k'}(t') \rangle = \delta(k + k')\delta(t - t') \cdot 2[1 - \cos(k)]
\end{equation}

\subsubsection{Final Expression for the Structure Factor}

Using the above correlated noise to solve Eq.~(\ref{eq:first-order-langevin}), we obtain the structure factor in the long-time limit. Note that the structure factor can be calculated from fluctuations of the correlated Poisson fields\cite{bertrand2019}:
\begin{equation}\label{eq:ns-definition}
N_S(k) = \langle \rho_k \rho_{-k} \rangle = \rho_0 + \langle \delta\tilde{\rho}_k \, \delta\tilde{\rho}_{-k} \rangle
\end{equation}

\begin{remark}[Why can $N_S(k)$ be written as $\rho_0 + \langle \delta\tilde{\rho}_k \delta\tilde{\rho}_{-k} \rangle$?]
This is a standard identity of the Poisson representation. The physical occupation numbers $n_i = n_i^+ + n_i^-$ are Poisson variables with rates $\alpha_i, \beta_i$. For Poisson processes we have:
\[
\langle n_i^\pm n_j^\pm \rangle = \langle \alpha_i \alpha_j \rangle + \delta_{ij}\langle \alpha_i \rangle, \qquad
\langle n_i^+ n_j^- \rangle = \langle \alpha_i \beta_j \rangle
\]
Therefore, the correlation of the total particle number is:
\begin{align*}
\langle n_i n_j \rangle
&= \langle (n_i^+ + n_i^-)(n_j^+ + n_j^-) \rangle \\
&= \langle (\alpha_i + \beta_i)(\alpha_j + \beta_j) \rangle + \delta_{ij}\langle \alpha_i + \beta_i \rangle \\
&= \langle \tilde{\rho}_i \tilde{\rho}_j \rangle + \delta_{ij}\rho_0
\end{align*}
The second term $\delta_{ij}\rho_0$ is precisely the \textbf{self-correlation} contribution from Poisson statistics (fluctuations in the particle number at the same lattice site follow a Poisson distribution, with variance equal to the mean). Transforming to Fourier space and considering $k \neq 0$ modes (for which $\tilde{\rho}_k = \delta\tilde{\rho}_k$, because the uniform background $\rho_0$ contributes only at $k=0$), we obtain Eq.~(\ref{eq:ns-definition}).
\end{remark}

\begin{derivation}[Detailed calculation of the structure factor]
The solution of Eq.~(\ref{eq:first-order-langevin}) can be obtained via Fourier transform. Let the steady-state fluctuation $\delta\tilde{\rho}_k(t)$ be:
\[
\delta\tilde{\rho}_k(t) = ic \int_{-\infty}^{t} e^{-Dk^2(t-t')}\eta_k(t')\,\dd t'
\]

Therefore:
\begin{align}
\langle \delta\tilde{\rho}_k \, \delta\tilde{\rho}_{-k} \rangle &= -c^2 \int_{-\infty}^{t}\int_{-\infty}^{t} e^{-Dk^2(2t - t' - t'')} \langle \eta_k(t')\eta_{-k}(t'') \rangle \, \dd t' \, \dd t'' \nonumber\\
&= -c^2 \int_{-\infty}^{t} e^{-2Dk^2(t-t')} \cdot 2[1 - \cos(k)] \, \dd t' \nonumber\\
&= -c^2 \frac{2[1 - \cos(k)]}{2Dk^2} \nonumber\\
&= -\frac{c^2[1 - \cos(k)]}{Dk^2}\nonumber
\end{align}
\end{derivation}

The final result is:
\begin{resultbox}[Structure Factor (Core Result)]
\begin{equation}\label{eq:structure-factor-final}
\boxed{N_S(k) = \rho_0 - \frac{c^2[1 - \cos(k)]}{Dk^2} = \rho_0 - \frac{c^2}{2D}\left(1 - \frac{k^2}{12} + \ldots\right)}
\end{equation}
\end{resultbox}

\subsubsection{Physical Analysis}

Comparing $S_0 = \lim_{k\to 0^+} N_S(k) = \rho_0 - c^2/(2D)$ with the Poisson value $S|_{r_1=0} = \rho_0$:
\begin{equation}\label{eq:reduction-ratio}
\frac{S_0}{\rho_0} = 1 - \frac{c^2}{2D\rho_0} = 1 - \frac{r_1\rho_0}{r_0 + r_1\rho_0}
\end{equation}

In the interaction-dominated limit $r_1\rho_0 \gg r_0$:
\begin{equation}\label{eq:limit-result}
\boxed{\frac{S_0}{\rho_0} \to \frac{1}{2}}
\end{equation}

This means that long-wavelength density fluctuations are suppressed to one half of the Poisson value!

\begin{notebox}[Key Physical Conclusions]
\begin{enumerate}
    \item When $r_1\rho_0 \gg r_0$ (strong antialigning interaction), $S_0/\rho_0 \to 1/2$
    \item Introducing spontaneous conversion $r_0$ is equivalent to adding orientational noise, supporting the early assertion that ``dynamically established correlations are more pronounced in deterministic models''
    \item Lower values (including hyperuniform $S_0 \to 0$) may be achieved in models with strict center-of-mass or polarization conservation
    \item In the $k \to 0$ limit, the structure factor approaches a \textbf{finite constant} rather than zero, so the system is \textbf{not strictly hyperuniform}, but exhibits significant fluctuation suppression
\end{enumerate}
\end{notebox}

\subsection{Numerical Verification (One-Dimensional Model)}\label{subsec:1d-numerics}

The paper also reports numerical simulation results for the one-dimensional model. The stochastic process was simulated directly using the Gillespie algorithm\cite{boltz2025}, with simulation parameters $\lambda = 1$, $r_1 = 1/100$, and variable system size $L$.

\begin{itemize}
    \item Simulation results are qualitatively consistent with the predictions of Eq.~(\ref{eq:structure-factor-final})
    \item For sufficiently large system sizes, suppression of long-wavelength density fluctuations is observed
    \item Note that in finite systems, there is a difference between Binomial and Poisson statistics, which requires correction
\end{itemize}

\begin{remark}[Finite-size effects and spurious hyperuniformity]
If the offset $S_0$ is ignored or unknown in the small-$k$ regime, it may lead to erroneous inferences---namely, inferring a spurious power-law behavior. In the case of Fig.~2, this could lead to inferring an ``anomalous'' scaling of $k^{\approx 0.25}$. This reminds us that inferring hyperuniformity from numerical data without a primary theoretical basis is risky.
\end{remark}
